\documentclass{article}
\usepackage[margin=1.5cm]{geometry}
\usepackage{amsmath}
\usepackage{enumitem}

\setlist{nosep,leftmargin=*}

\begin{document}
\twocolumn

\title{Chapter 8 In-Class Exercise: Searching Strings in a Patient Database}
\author{Prof. Jordan C. Hanson}
\date{September 21, 2026}

\maketitle

\section{Starting Data}

\begin{enumerate}

\item Chapter 8 string operations allow us to make a database
search more flexible.  Instead of requiring an exact patient ID,
we can search for names and partial names.

Begin with the patient database

\begin{verbatim}
patients = [
 [1042, "Maria Garcia", 37],
 [1057, "Daniel Kim", 52],
 [1081, "Aisha Patel", 29],
 [1094, "Elena Garcia", 44],
 [1103, "Marcus Johnson", 61]
]
\end{verbatim}

Each record contains

\begin{verbatim}
[patient_id, name, age]
\end{verbatim}

\end{enumerate}


\section{Exact String Match}

\begin{enumerate}

\item Write a loop that prints the complete patient record whose
name is exactly

\begin{verbatim}
"Maria Garcia"
\end{verbatim}

Use the string comparison operator \texttt{==}.  A useful structure is

\begin{verbatim}
for patient in patients:
    if patient[1] == "Maria Garcia":
        print(patient)
\end{verbatim}

How many patient records should be printed?

\end{enumerate}

\section{Substring Match}

\begin{enumerate}

\item Modify the search so that every patient whose name contains \verb+"Garcia"+ is printed.  Use Python's \texttt{in} operator:

\begin{verbatim}
if "Garcia" in patient[1]:
    print(patient)
\end{verbatim}

\begin{itemize}

\item How many records are printed?

\item Why does this search find both \texttt{"Maria Garcia"} and \texttt{"Elena Garcia"}?

\end{itemize}

\end{enumerate}


\section{Case-Insensitive Search}

\begin{enumerate}

\item A user may type ''garcia'', ''Garcia'', or ''GARCIA''.  Use the string method \verb+.lower()+ so that capitalization does not affect the search. Try

\begin{verbatim}
search = "GARCIA"
for patient in patients:
    if search.lower() in patient[1].lower():
		print(patient)
\end{verbatim}

Why must both \texttt{search} and the patient name be converted to lowercase?

\end{enumerate}

\section{User-Controlled Search}

\begin{enumerate}

\item Ask the user for a complete or partial patient name:

\begin{verbatim}
search = input("Enter a patient name or partial name: ")
\end{verbatim}

Write a loop that searches all patient names and prints every matching record.  Your search must be case-insensitive. Test your program using

\begin{verbatim}
Garcia
garcia
Daniel
PATEL
\end{verbatim}

\end{enumerate}


\section{Challenge: Count Matches}

\begin{enumerate}

\item Add a counter before the loop:

\begin{verbatim}
matches = 0
\end{verbatim}

Whenever a matching patient is found, increase the counter:

\begin{verbatim}
matches = matches + 1
\end{verbatim}

After the search finishes, use

\begin{verbatim}
if matches == 0:
    print("No matching patients "
          "were found.")
else:
    print("Matches:", matches)
\end{verbatim}

Test the program with one search that produces multiple matches
and another search that produces no matches.

\vspace{0.15cm}

\textbf{Key Chapter 8 ideas:}
string equality with \texttt{==}, substring membership with
\texttt{in}, and string methods such as \texttt{.lower()}.

\end{enumerate}

\end{document}